% Artigo P2 de PDI --- template Elsevier (elsarticle), formato Results in Engineering.
% Compilar no Overleaf (pdfLaTeX + BibTeX) ou localmente: latexmk -pdf main.tex
\documentclass[preprint,12pt,a4paper]{elsarticle}

\usepackage[utf8]{inputenc}
\usepackage[T1]{fontenc}
\usepackage[brazil]{babel}
\usepackage{amsmath,amssymb,amsthm}
\usepackage{booktabs,longtable,multirow,array}
\usepackage{graphicx}
\usepackage{xcolor}
\usepackage{siunitx}
\sisetup{output-decimal-marker={,}}
\usepackage{tikz}
\usetikzlibrary{positioning,arrows.meta,shapes.geometric,fit,backgrounds}
\usepackage[hidelinks]{hyperref}
\usepackage{float}
\usepackage{ifthen}

% ---------------------------------------------------------------------------
% Origem dos resultados: "piloto" (CPU, reduzido) ou "exp" (experimento completo em GPU).
% Depois de rodar o experimento completo, trocar para {exp} e recompilar.
\newcommand{\resdir}{piloto}
\graphicspath{{figs/}}
\newcommand{\figres}[1]{figs/\resdir/#1}
\newcommand{\tabres}[1]{\input{tabelas/\resdir/#1}}
% Marcação de texto a revisar com os números do experimento completo.
\newcommand{\revisar}[1]{\textcolor{orange!70!black}{#1}}
% ---------------------------------------------------------------------------

\journal{Results in Engineering}

\begin{document}

\begin{frontmatter}

\title{Segmentação de padrões de Turing do sistema de Gray-Scott com U-Net, SegFormer e
Mask2Former: efeito de dez pré-processamentos e cinco pós-processamentos}

\author[ufma]{Equipe Gray-Scott}
\affiliation[ufma]{organization={Universidade Federal do Maranhão (UFMA)},
            city={São Luís}, state={MA}, country={Brasil}}

\begin{abstract}
Investigamos onde estão as estruturas de interesse num sistema de reação-difusão e como elas
evoluem no tempo, segmentando os padrões de Turing do conjunto \texttt{gray\_scott\_reaction\_diffusion}
do The Well (seis regimes $(f,k)$, campos $A$ e $B$ em $128\times128$ com contorno periódico). A
rede recebe só o campo $A$, e o rótulo vem de um critério físico sobre $B$ ($B>\tau$, com $\tau$
de Otsu ajustado no treino), o que evita vazamento de rótulo. Seguindo o protocolo comum da turma,
três arquiteturas (U-Net, SegFormer e Mask2Former) são treinadas por 20 épocas, com o checkpoint
escolhido pelo Dice de validação. Elas são avaliadas em quatro condições: sem pré e sem pós,
somente pré (P1--P10), somente pós (Q1--Q5) e pré e pós (150 combinações). Usamos todas as 1200
trajetórias da divisão oficial e reportamos IoU, Dice, Precision e Recall com IC de 95\% por
bootstrap sobre trajetórias, além de métricas temporais e de custo, por regime e no agregado, no
formato padronizado do benchmark entre sistemas físicos.
\end{abstract}

\begin{keyword}
Gray-Scott \sep reação-difusão \sep padrões de Turing \sep segmentação semântica \sep U-Net \sep
SegFormer \sep Mask2Former \sep pré-processamento \sep pós-processamento \sep The Well
\end{keyword}

\end{frontmatter}

\section{Introdução}
\label{sec:intro}

Sistemas de reação-difusão produzem, a partir de um estado quase homogêneo, estruturas espaciais
organizadas (pontos, listras, labirintos, espirais e ``planadores'') cuja forma depende de poucos
parâmetros \cite{turing1952,pearson1993}. A pergunta que orienta este trabalho é a mesma para
todos os sistemas físicos estudados pela turma: \emph{onde estão as estruturas de interesse e como
elas evoluem ao longo do tempo?} Responder a ela exige uma segmentação por pixel, e não apenas a
classificação do regime de cada simulação.

O sistema estudado pela equipe é o \texttt{gray\_scott\_reaction\_diffusion} do The Well
\cite{ohana2024well}: 1200 trajetórias do modelo de Gray-Scott \cite{grayscott1984,pearson1993}
em seis regimes de parâmetros, com 1001 instantes de $128\times128$ pixels cada. O conjunto não
traz máscaras de segmentação, e o ground truth (GT) é derivado da concentração $B$ por um critério
físico documentado e congelado antes da comparação (Seção~\ref{sec:gt}).

O experimento segue o protocolo comum da turma, que permite juntar os resultados de 11 sistemas
físicos 2D num único benchmark:
\begin{itemize}
  \item resolução com maior eixo de no máximo 128 px e sem ampliação;
  \item no máximo 25 GB de dados por sistema;
  \item três arquiteturas: U-Net \cite{ronneberger2015}, SegFormer \cite{xie2021segformer} e
    Mask2Former \cite{cheng2022mask2former};
  \item 20 épocas, com o checkpoint da melhor época de validação;
  \item dez pré-processamentos (P1--P10) e cinco pós-processamentos (Q1--Q5), cada um testado
    isoladamente contra o baseline.
\end{itemize}

São avaliadas quatro condições:
\begin{enumerate}
  \item sem pré e sem pós (referência);
  \item somente pré (efeito isolado de P1--P10);
  \item somente pós (efeito isolado de Q1--Q5);
  \item pré e pós (as $10\times3\times5 = 150$ combinações).
\end{enumerate}

As contribuições são:
\begin{itemize}
  \item um protocolo reprodutível de segmentação do Gray-Scott, com GT físico, entrada restrita a
    $A$ para evitar vazamento de rótulo, divisão oficial por trajetória e operações que respeitam o
    contorno periódico;
  \item a análise de \emph{por que} cada pré e pós-processamento melhora ou piora a segmentação,
    relacionando o efeito ao mecanismo de cada arquitetura (convolução com \emph{skip connections},
    atenção hierárquica em 1/4 da resolução e classificação de máscaras por consultas) e à física
    de cada regime;
  \item tabelas de resultados no esquema padronizado do protocolo (CSV e \LaTeX), prontas para a
    junção com os demais sistemas no artigo de benchmark.
\end{itemize}

\section{Onde o fluido se aplica}
\label{sec:aplicacoes}

O Gray-Scott é um sistema de reação-difusão, não um escoamento no sentido de Navier-Stokes. Ele
entra no conjunto de ``fluidos'' do The Well como um meio contínuo descrito por EDPs de
transporte (difusão) acopladas a fontes não lineares (reação). As aplicações do modelo e da
segmentação dos seus padrões são as seguintes.

\paragraph{Química} O modelo descreve a autocatálise cúbica $A + 2B \to 3B$, $B \to C$ num
reator alimentado continuamente \cite{grayscott1984}. Os padrões previstos por Pearson
\cite{pearson1993} foram observados experimentalmente no sistema ferrocianeto-iodato-sulfito em gel
\cite{lee1993}. Segmentar as manchas permite medir número, área, velocidade e divisão de
estruturas, que são as grandezas comparadas entre experimento e simulação.

\paragraph{Biologia e ecologia} Mecanismos de Turing explicam padrões de pigmentação (zebras,
peixes-zebra, conchas) e o espaçamento de estruturas no desenvolvimento \cite{kondo2010,murray2003},
assim como os padrões de vegetação em regiões semiáridas \cite{klausmeier1999}. Nesses domínios, a
máscara segmentada é a observável primária: fração ocupada, comprimento de onda e conectividade
dos padrões.

\paragraph{Computação gráfica: o pipeline gráfico} Reação-difusão é uma técnica clássica de
síntese procedural de texturas \cite{turk1991,witkin1991}. No pipeline gráfico, o sistema é
integrado no espaço de textura (ou diretamente na malha \cite{turk1991}), e o campo resultante é
mapeado em cor, rugosidade ou deslocamento na etapa de \emph{shading}. A segmentação de $B$ é o
que transforma o campo contínuo em regiões (por exemplo, manchas de pelagem com bordas nítidas e
máscaras de material). Ela também é usada para controlar a evolução do padrão quadro a quadro em
animações, o que corresponde ao filtro temporal Q10. Implementações em GPU (\emph{compute} ou
\emph{fragment shaders} com textura em \emph{ping-pong}) usam o mesmo esquema de diferenças com
contorno periódico (\emph{wrap}) adotado aqui.

\paragraph{Aprendizado de máquina científico} O The Well usa o Gray-Scott como benchmark de
previsão de campos (modelos substitutos) \cite{ohana2024well}. Uma segmentação confiável das
estruturas fornece métricas físicas interpretáveis para validar esses modelos, como a preservação
do número de manchas e a coerência temporal das frentes. Essas métricas complementam erros
pixel a pixel como o VRMSE.

\paragraph{Por que segmentar como imagem} Os campos do The Well estão em grades regulares, e o
protocolo da turma limita todos a no máximo 128 px no maior eixo. Assim, o mesmo conjunto de
arquiteturas de segmentação semântica de imagens pode ser aplicado a todos os sistemas físicos. As
três redes cobrem as três famílias dominantes: convolucional com \emph{skip connections} (U-Net),
\emph{Transformer} hierárquico com decodificador leve (SegFormer) e classificação de máscaras por
consultas (Mask2Former). A comparação entre sistemas físicos mede, portanto, como as propriedades de
cada campo (escala das estruturas, contraste e dinâmica) interagem com o viés indutivo de cada
família.

\section{Fundamentação teórica}
\label{sec:teoria}

\subsection{O modelo de Gray-Scott}
\label{sec:grayscott}

O modelo descreve duas espécies químicas com concentrações $A(\mathbf{x},t)$ (substrato) e
$B(\mathbf{x},t)$ (autocatalisador) num domínio $\Omega=[-1,1]^2$ com contorno periódico
\cite{grayscott1984,pearson1993}:
\begin{align}
  \frac{\partial A}{\partial t} &= \delta_A \nabla^2 A - A B^2 + f\,(1 - A), \label{eq:gsA}\\
  \frac{\partial B}{\partial t} &= \delta_B \nabla^2 B + A B^2 - (f + k)\,B. \label{eq:gsB}
\end{align}
O termo $AB^2$ é a reação autocatalítica $A + 2B \to 3B$: consome $A$ e produz $B$. O termo
$f(1-A)$ é a \emph{alimentação} (\emph{feed}), que repõe $A$ em direção ao reservatório $A=1$ à
taxa $f$. O termo $-(f+k)B$ é a \emph{remoção}: $B$ sai do sistema por arraste ($f$) e decai para
um produto inerte $C$ ($k$, \emph{kill}). No The Well, $\delta_A = 2\cdot10^{-5}$ e
$\delta_B = 10^{-5}$ em todas as simulações \cite{ohana2024well}.
A documentação oficial do conjunto escreve $-AB^2$ na Eq.~\eqref{eq:gsB}; o sinal correto,
usado pelo gerador dos dados, é $+AB^2$.

\paragraph{Estados homogêneos} Igualando os termos de reação a zero, o estado trivial
$(A,B)=(1,0)$ existe para todo $(f,k)$. Para $B\neq0$, a Eq.~\eqref{eq:gsB} dá $AB = f+k$, e a
substituição na Eq.~\eqref{eq:gsA} leva a
\begin{equation}
  (f+k)\,B^2 - f\,B + f\,(f+k) = 0 \;\Rightarrow\;
  B_\pm = \frac{f \pm \sqrt{f^2 - 4 f (f+k)^2}}{2(f+k)}, \qquad A_\pm = \frac{f+k}{B_\pm},
\end{equation}
que existem quando $f \geq 4(f+k)^2$. O Jacobiano no estado trivial é
$\mathrm{diag}(-f, -(f+k))$, com autovalores negativos: o fundo $(1,0)$ é \emph{linearmente
estável}. Os padrões não surgem de uma instabilidade do fundo, mas de perturbações finitas, que
são as condições iniciais (séries de Fourier e aglomerados gaussianos) do The Well. Isso explica
por que algumas trajetórias decaem para o fundo homogêneo (as trajetórias estacionárias de Gliders
e Spirals).

\paragraph{Instabilidade de Turing} Seja $\mathbf{J}$ o Jacobiano num estado homogêneo não
trivial e $\mathbf{D}=\mathrm{diag}(\delta_A,\delta_B)$. Uma perturbação
$\propto e^{\lambda t + i\mathbf{q}\cdot\mathbf{x}}$ cresce se algum autovalor de
$\mathbf{J}-q^2\mathbf{D}$ tiver parte real positiva. Com
$J_{11} = -B^2 - f$, $J_{12} = -2AB$, $J_{21} = B^2$ e $J_{22} = 2AB - (f+k)$, a instabilidade
induzida por difusão ocorre quando o estado é estável sem difusão
($\mathrm{tr}\,\mathbf{J}<0$, $\det\mathbf{J}>0$) e
\begin{equation}
  h(q^2) = \delta_A\delta_B\, q^4 - (\delta_A J_{22} + \delta_B J_{11})\, q^2 + \det\mathbf{J} < 0
  \quad\text{para algum } q,
\end{equation}
isto é, $\delta_A J_{22} + \delta_B J_{11} > 2\sqrt{\delta_A\delta_B\det\mathbf{J}}$
\cite{turing1952,murray2003,cross2009}. O número de onda crítico
$q_c^2 = (\delta_A J_{22} + \delta_B J_{11})/(2\delta_A\delta_B)$ fixa o comprimento de onda
característico $2\pi/q_c$ do padrão. Ele depende de $(f,k)$ e de $\sqrt{\delta}$ e, nesta
discretização, corresponde a estruturas de 4 a 8 pixels de largura. Esse é o motivo pelo qual a
resolução não pode ser reduzida sem destruir o padrão (Seção~\ref{sec:resolucao}).

\paragraph{Os seis regimes} Pearson \cite{pearson1993} classificou o plano $(f,k)$ por
morfologia. O The Well usa seis pontos (Tabela~\ref{tab:regimes_fk}): \emph{Gliders}
(estruturas móveis que se propagam e se aniquilam), \emph{Bubbles} (bolhas de fundo dentro de uma
fase de $B$), \emph{Maze} (labirinto de listras quase todo conexo), \emph{Worms} (vermes
alongados), \emph{Spirals} (ondas espirais em movimento permanente) e \emph{Spots} (manchas que se
replicam por divisão até preencher o domínio).

\begin{table}[t]\centering\small
\caption{Regimes do conjunto e características medidas na auditoria (Etapa 1): fração de pixels de
padrão ($B>\tau$, validação) e instante aproximado de fim do transiente.}
\label{tab:regimes_fk}
\begin{tabular}{lcccl}
\toprule
Regime & $f$ & $k$ & Pixels de padrão & Dinâmica após o transiente \\
\midrule
Gliders & 0,014 & 0,054 & 5\%  & móvel; 44\% das trajetórias decaem ao fundo \\
Bubbles & 0,098 & 0,057 & 28\% & evolução lenta \\
Maze    & 0,029 & 0,057 & 52\% & quase estático após $\sim$1500 s \\
Worms   & 0,058 & 0,065 & 25\% & evolução lenta \\
Spirals & 0,018 & 0,051 & 27\% & móvel permanente \\
Spots   & 0,030 & 0,062 & 26\% & quase estático após $\sim$3500 s \\
\bottomrule
\end{tabular}
\end{table}

\paragraph{Simulação} Os dados foram gerados por um método espectral de Fourier com integração
temporal exponencial de quarta ordem (ETDRK4 \cite{kassam2005}), via Chebfun \cite{chebfun2014},
com passo de 1 s e campos salvos a cada 10 s, de $t=0$ a $t=10\,000$ s. Por isso os dados são
\emph{praticamente livres de ruído}. Esse fato é central na interpretação dos filtros P6--P10
(Seção~\ref{sec:discussao}).

\subsection{O frame como nuvem de pontos}
\label{sec:nuvem}

Um frame é a função discreta $a:\mathcal{G}\to\mathbb{R}$ sobre a grade
$\mathcal{G}=\{(x_i,y_i)\}_{i=1}^{N}$. Na representação por nuvem de pontos ele é o conjunto
\begin{equation}
  \mathcal{S} = \{\,\mathbf{p}_i = (x_i, y_i, a_i)\,\}_{i=1}^{N}, \qquad a_i = \mathcal{P}(A)(x_i,y_i),
\end{equation}
em que $\mathcal{P}$ é o pré-processamento. A tarefa é aprender $g_\theta:\mathcal{S}\mapsto
\{\hat{y}_i\}$ com $\hat{y}_i\in[0,1]$, \emph{equivariante a permutações}: reordenar os pontos
reordena as saídas da mesma forma. Geometricamente, $\mathcal{S}$ é uma superfície de altura
$z=a(x,y)$ amostrada em $\mathbb{R}^3$. As estruturas do Gray-Scott são vales de $A$ (onde $B$ é
alto), e segmentar é separar os pontos dos vales dos pontos do platô $a\approx1$.

\subsection{PointNet}
PointNet \cite{qi2017pointnet} aproxima funções de conjunto pela composição
\begin{equation}
  F(\mathcal{S}) \approx \gamma\Big(\max_{i=1..N} h(\mathbf{p}_i)\Big),
\end{equation}
em que $h$ é um MLP compartilhado por ponto, $\max$ é o \emph{max-pooling} por canal (função
simétrica) e $\gamma$ é um MLP. O teorema de aproximação universal do artigo mostra que qualquer
função contínua de conjuntos (na métrica de Hausdorff) é aproximada por essa forma, com o mesmo
princípio dos \emph{Deep Sets} \cite{zaheer2017deepsets}. Na segmentação, o descritor global
$\mathbf{g}=\max_i h(\mathbf{p}_i)$ é concatenado à característica local de cada ponto, e um MLP
prediz $\hat{y}_i = \sigma(\phi([h_1(\mathbf{p}_i);\mathbf{g}]))$. Cada ponto é classificado só a
partir de si mesmo e de um resumo global. \textbf{Não há vizinhança local}: a rede não ``vê''
gradientes nem formas, apenas o valor $a_i$, a posição e o contexto global. Usamos a versão sem
T-Nets, porque a nuvem é 2D com orientação canônica e as rotações de $90^\circ$ entram como aumento
de dados.

\subsection{PointNet++}
PointNet++ \cite{qi2017pointnet2} aplica PointNet de forma hierárquica em vizinhanças
(\emph{set abstraction}, SA). Em cada nível $\ell$:
\begin{enumerate}
  \item \emph{amostragem}: $M_\ell$ centroides por amostragem do ponto mais distante (FPS), que
    cobre o conjunto de forma aproximadamente uniforme;
  \item \emph{agrupamento}: os $k$ vizinhos mais próximos $\mathcal{N}(c)$ de cada centroide $c$;
  \item \emph{PointNet local}: $\mathbf{f}_c = \max_{j\in\mathcal{N}(c)} h_\ell\big([\mathbf{x}_j-\mathbf{x}_c;\ \mathbf{f}_j]\big)$.
\end{enumerate}
As coordenadas relativas $\mathbf{x}_j-\mathbf{x}_c$ tornam as características invariantes a
translações. Usamos três níveis ($4096\to1024\to256\to64$ pontos, $k=16$). Na decodificação
(\emph{feature propagation}), as características de um nível esparso são interpoladas para o nível
denso por inverso da distância com os 3 vizinhos,
\begin{equation}
  \mathbf{f}(\mathbf{x}) = \frac{\sum_{j=1}^{3} w_j \mathbf{f}_j}{\sum_{j=1}^{3} w_j},\qquad
  w_j = \frac{1}{\|\mathbf{x} - \mathbf{x}_j\|},
\end{equation}
e concatenadas às do mesmo nível do codificador (\emph{skip connections}), de modo análogo à U-Net.
O campo receptivo cresce com a profundidade e permite reconhecer formas (manchas, listras) e não
apenas valores.

\subsection{DGCNN}
O DGCNN \cite{wang2019dgcnn} define a convolução de arestas (EdgeConv) sobre um grafo $k$-NN,
\begin{equation}
  \mathbf{x}_i' = \max_{j:(i,j)\in\mathcal{E}} h_\Theta\big(\mathbf{x}_i,\ \mathbf{x}_j - \mathbf{x}_i\big),
\end{equation}
que combina a característica do ponto (informação global) com as diferenças em relação aos vizinhos
(informação local, análoga a um gradiente discreto). O grafo é \emph{dinâmico}: na primeira camada
os vizinhos são definidos pelas coordenadas $(x,y)$ (grafo espacial). Nas camadas seguintes, são
recalculados no espaço de características, de modo que pontos distantes mas semelhantes (por
exemplo, duas manchas do mesmo regime) passam a trocar informação. Três EdgeConv
($k=16$), uma projeção global de 1024 canais e um MLP por ponto formam a cabeça de segmentação.

\subsection{Pré-processamentos}
\label{sec:pre}
Os pré-processamentos se dividem em transformações numéricas pontuais (P1--P5, que mudam só os
valores de $a_i$) e filtros espaciais (P6--P10, que mudam a forma do campo)
\cite{gonzalez2018}. Estatísticas ($\mu,\sigma$, quantis $Q_p$) são calculadas \textbf{só no
treino} e aplicadas sem recalibração à validação e ao teste. Todos os filtros usam extensão
periódica.
\begin{description}
  \item[P1 Z-score:] $a = (A-\mu)/\sigma$.
  \item[P2 Min-Max:] $a = (A-\min A)/(\max A-\min A)$.
  \item[P3 Robusta:] $a = (A - Q_{50})/(Q_{75}-Q_{25})$.
  \item[P4 Recorte por percentis:] $a = \min(\max(A, Q_{1}), Q_{99})$, sem reescala, para isolar o
    efeito do recorte.
  \item[P5 Logarítmica:] $a = \log(\max(A,0)+\varepsilon)$, $\varepsilon = 10^{-3}$, que expande os
    valores baixos de $A$ (o interior das estruturas).
  \item[P6 Gaussiano:] $a = G_\sigma * A$, $\sigma = 1$ px.
  \item[P7 Mediana:] $a(\mathbf{x}) = \mathrm{med}\{A(\mathbf{y}) : \mathbf{y}\in W_{3\times3}(\mathbf{x})\}$.
  \item[P8 Bilateral] \cite{tomasi1998}: $a(\mathbf{x}) = \frac{1}{Z}\sum_{\mathbf{y}} e^{-\|\mathbf{x}-\mathbf{y}\|^2/2\sigma_s^2}\, e^{-(A(\mathbf{x})-A(\mathbf{y}))^2/2\sigma_r^2}A(\mathbf{y})$,
    com $\sigma_s=1$ px, $\sigma_r=0{,}05$ e janela $5\times5$.
  \item[P9 Difusão anisotrópica] \cite{perona1990}: $\partial_\tau a = \nabla\cdot\big(g(\|\nabla a\|)\nabla a\big)$,
    $g(s) = e^{-(s/\kappa)^2}$, $\kappa=0{,}05$, 10 iterações explícitas com $\Delta\tau=0{,}2$.
  \item[P10 Wavelets] \cite{donoho1994,chang2000}: decomposição Daubechies-2 em 2 níveis com
    extensão periódica e limiarização suave BayesShrink,
    $\tau_s = \hat\sigma^2/\hat\sigma_{X,s}$, com $\hat\sigma = \mathrm{med}(|d_{HH_1}|)/0{,}6745$.
\end{description}

\subsection{Pós-processamentos}
\label{sec:pos}
A rede produz a probabilidade $P(\mathbf{x})$ do ponto pertencer ao padrão. Sem pós, a máscara é
$\hat{M} = \mathbb{1}[P>0{,}5]$ (Q0). Os parâmetros de cada método são escolhidos na validação
(Tabela~\ref{tab:grades}). A morfologia \cite{serra1982,soille2003} e a rotulação de componentes
são periódicas: um componente que toca uma borda continua na borda oposta.
\begin{description}
  \item[Q1 Limiarização:] $\hat{M} = \mathbb{1}[P > \tau^*]$, com $\tau^*$ que maximiza o Dice de
    validação.
  \item[Q2 Abertura:] $\hat{M}\circ S = (\hat{M}\ominus S)\oplus S$, que remove saliências e objetos
    menores que o elemento estruturante $S$.
  \item[Q3 Fechamento:] $\hat{M}\bullet S = (\hat{M}\oplus S)\ominus S$, que fecha lacunas.
  \item[Q4 Abertura + fechamento:] $(\hat{M}\circ S)\bullet S$.
  \item[Q5 Remoção de componentes pequenos:] remove os componentes $C$ com $|C| < a_{\min}$.
  \item[Q6 Preenchimento de buracos:] no toro não há ``exterior'', então um buraco é definido como um
    componente do fundo com área $\leq a_{\max}$, que é incorporado ao padrão.
  \item[Q7 Seleção de componentes:] mantém os componentes com probabilidade média
    $\bar{P}_C \geq p_{\min}$. O critério ``maior componente'' não se aplica, porque Spots tem
    $\sim$200 componentes por frame.
  \item[Q8 Watershed] \cite{vincent1991}: inundação de $\|\nabla P\|$ (Sobel) a partir de marcadores
    de padrão ($P>p_{\text{alto}}$) e de fundo ($P<p_{\text{baixo}}$). A fronteira é posta na
    crista do gradiente, e não no nível 0,5.
  \item[Q9 CRF] \cite{lafferty2001,krahenbuhl2011}: minimiza
    $E(\mathbf{y}) = \sum_i -\log P(y_i) + \sum_{i<j}\mu(y_i,y_j)\,k(\mathbf{f}_i,\mathbf{f}_j)$
    com o modelo de Potts $\mu$ e o kernel
    $k = w_a e^{-\frac{\|\mathbf{x}_i-\mathbf{x}_j\|^2}{2\theta_\alpha^2}-\frac{(A_i-A_j)^2}{2\theta_\beta^2}} + w_s e^{-\frac{\|\mathbf{x}_i-\mathbf{x}_j\|^2}{2\theta_\gamma^2}}$,
    por 5 iterações de campo médio numa janela periódica $7\times7$.
  \item[Q10 Filtragem temporal:] $\tilde{P}_t = \frac{1}{w}\sum_{s=-\lfloor w/2\rfloor}^{\lfloor w/2\rfloor} P_{t+s}$,
    seguida de $\mathbb{1}[\tilde{P}_t>0{,}5]$, ao longo dos frames de cada trajetória (passo de 200 s).
\end{description}

\begin{table}[t]\centering\small
\caption{Grades de parâmetros dos pós-processamentos (escolha pelo Dice médio de validação).}
\label{tab:grades}
\begin{tabular}{ll}
\toprule
Método & Grade \\
\midrule
Q1 & $\tau\in\{0{,}05, 0{,}10, \dots, 0{,}95\}$ \\
Q2, Q3, Q4 & $S\in\{\text{cruz } 3\times3,\ \text{quadrado } 3\times3\}$ \\
Q5 & $a_{\min}\in\{2,4,8,16,32\}$ px \\
Q6 & $a_{\max}\in\{2,4,8,16,32\}$ px \\
Q7 & $p_{\min}\in\{0{,}6; 0{,}7; 0{,}8; 0{,}9\}$ \\
Q8 & $p_{\text{alto}}\in\{0{,}7; 0{,}8; 0{,}9\}$, $p_{\text{baixo}}\in\{0{,}1; 0{,}2; 0{,}3\}$ \\
Q9 & $w_a\in\{1,3\}$, $\theta_\beta\in\{0{,}05; 0{,}1\}$, $w_s=1$, $\theta_\alpha=3$, $\theta_\gamma=1$ \\
Q10 & $w\in\{3,5\}$ frames \\
\bottomrule
\end{tabular}
\end{table}

\subsection{Métricas}
\label{sec:metricas}
Com $\mathrm{TP}$, $\mathrm{FP}$ e $\mathrm{FN}$ contados por pixel (ponto):
\begin{gather}
  \mathrm{IoU} = \frac{\mathrm{TP}}{\mathrm{TP}+\mathrm{FP}+\mathrm{FN}} \;\text{\cite{jaccard1912}},\qquad
  \mathrm{Dice} = \frac{2\,\mathrm{TP}}{2\,\mathrm{TP}+\mathrm{FP}+\mathrm{FN}} \;\text{\cite{dice1945}},\\
  \mathrm{Precision} = \frac{\mathrm{TP}}{\mathrm{TP}+\mathrm{FP}},\qquad
  \mathrm{Recall} = \frac{\mathrm{TP}}{\mathrm{TP}+\mathrm{FN}}.
\end{gather}
Dice e IoU são monotonamente relacionados, $\mathrm{Dice} = 2\,\mathrm{IoU}/(1+\mathrm{IoU})$, e
portanto ordenam os métodos da mesma forma numa única máscara. O Dice é a média harmônica de
Precision e Recall. A Precision penaliza falso padrão (por exemplo, estruturas ``inventadas'' pela
rede) e a Recall penaliza estruturas perdidas (por exemplo, Gliders pequenos apagados por uma
abertura). Seguindo a convenção G4, TP, FP e FN são somados nos frames de cada trajetória, e a
métrica é calculada por trajetória (vazio contra vazio vale 1).

\paragraph{Métricas temporais} Para uma trajetória com máscaras $\hat{M}_t$ e GT $M_t$ em frames
consecutivos:
\begin{itemize}
  \item erro de variação de área:
    $E_{\text{área}} = \overline{|\Delta|\hat{M}_t|/|\Omega| - \Delta|M_t|/|\Omega||}$;
  \item erro de centroide:
    $E_{\text{centr.}} = \overline{d_{\mathbb{T}}(\hat{\mathbf{c}}_t, \mathbf{c}_t)}$, com o
    centroide circular \cite{bai2008} (média angular em cada eixo, adequada ao toro) e a distância
    periódica $d_{\mathbb{T}}$;
  \item erro de consistência:
    $E_{\text{cons.}} = \overline{|\mathrm{IoU}(\hat{M}_t,\hat{M}_{t+1}) - \mathrm{IoU}(M_t,M_{t+1})|}$.
\end{itemize}
Comparar com o GT, e não exigir apenas que a máscara mude pouco, evita premiar uma máscara
``congelada'' em regimes móveis (Spirals e Gliders), em que a mudança entre frames é física.

\paragraph{Incerteza} A unidade estatística é a trajetória, porque frames da mesma trajetória são
correlacionados. Reportamos a média sobre trajetórias e a semi-amplitude do intervalo de confiança
de 95\% por bootstrap percentil com 2000 reamostragens \cite{efron1993}.

\paragraph{Custo} A latência é o tempo médio por frame $128\times128$ (pré, mais rede com todos os
blocos do frame num lote, mais pós), em ms. A memória é o pico de memória de GPU alocada na
inferência, em MB.

\section{Metodologia}
\label{sec:metodo}

\subsection{Pipeline}
A Figura~\ref{fig:pipeline} resume o pipeline. Os campos são lidos por requisições de intervalo
de bytes diretamente dos HDF5 do The Well (os arquivos são contíguos e sem compressão), sem
baixar os 154 GB do conjunto.

\begin{figure}[t]
\centering
\resizebox{\textwidth}{!}{%
\begin{tikzpicture}[
  node distance=4mm and 5mm,
  box/.style={draw, rounded corners=2pt, align=center, font=\scriptsize, minimum height=10mm,
              text width=21mm, fill=#1},
  box/.default=gray!8,
  arr/.style={-{Stealth[length=2mm]}, thick, gray!70!black}]
  \node[box=blue!8] (well) {The Well\\HDF5 $A,B$\\$1001\times128\times128$};
  \node[box, right=of well] (split) {Divisão oficial\\por trajetória\\(treino/val./teste)};
  \node[box, right=of split] (sub) {Seleção de frames\\exclui $t=0$ e\\padrão morto};
  \node[box=orange!12, above right=2mm and 5mm of sub] (gt) {GT: $M=\mathbb{1}[B>\tau]$\\$\tau$ = Otsu (treino)};
  \node[box=green!10, below right=2mm and 5mm of sub] (pre) {Pré P0--P10\\sobre $A$\\(ajuste no treino)};
  \node[box=green!10, right=18mm of sub] (pc) {Nuvem de pontos\\$(x,y,a)$, blocos\\$64\times64$ periódicos};
  \node[box=violet!10, right=of pc] (net) {PointNet\\PointNet++\\DGCNN\\(20 épocas)};
  \node[box, right=of net] (prob) {Probabilidade\\por ponto $\to$\\grade (média\\de blocos)};
  \node[box=green!10, right=of prob] (pos) {Pós Q0--Q10\\(parâmetros\\na validação)};
  \node[box=blue!8, right=of pos] (met) {IoU, Dice, P, R\\temporais, custo\\IC 95\% $\to$ CSV};
  \draw[arr] (well) -- (split);
  \draw[arr] (split) -- (sub);
  \draw[arr] (sub) |- (gt);
  \draw[arr] (sub) |- (pre);
  \draw[arr] (pre) -| (pc);
  \draw[arr] (gt) -| (pc);
  \draw[arr] (pc) -- (net);
  \draw[arr] (net) -- (prob);
  \draw[arr] (prob) -- (pos);
  \draw[arr] (pos) -- (met);
  \draw[arr, dashed] (gt.east) -- ++(78mm,0) -| (met.north);
\end{tikzpicture}}
\caption{Pipeline do experimento. Os blocos em verde são os fatores variados (pré, pós); o GT
(laranja) é fixo e independente do modelo.}
\label{fig:pipeline}
\end{figure}

\subsection{Dados e divisão}
Usamos a divisão oficial do The Well, por trajetória, de modo que nenhuma simulação aparece em
mais de uma partição. Os seis regimes estão nas três partições (160/20/20 trajetórias por regime).
Por custo computacional, selecionamos de forma determinística trajetórias igualmente espaçadas nos
índices, o que cobre os dois tipos de condição inicial, que ocupam metades distintas de cada
arquivo (Tabela~\ref{tab:dados}).

\begin{table}[t]\centering\small
\caption{Subconjuntos usados (por regime). Índice de frame $n$ corresponde a $t = 10\,n$ s.}
\label{tab:dados}
\begin{tabular}{lcl}
\toprule
Partição & Trajetórias & Frames \\
\midrule
Treino & 24 de 160 & $n\in\{1,2,4,6,8,10,15,20,30,\dots,100\}$ (16; denso no transiente) \\
Validação & 5 de 20 & $n = 20, 40, \dots, 1000$ (50) \\
Teste & 10 de 20 & $n = 20, 40, \dots, 1000$ (50) \\
Tabela visual & trajetória 0 do teste & $n\in\{0, 40, 500, 1000\}$ \\
\bottomrule
\end{tabular}
\end{table}

\subsection{Ground truth}
\label{sec:gt}
O professor pediu que o GT fosse ``uma imagem em um frame médio que caracterize o fluido''. A ficha
da Etapa 2 comparou três interpretações. Uma máscara fixa de um instante intermediário e a média
temporal limiarizada falham nos regimes móveis: o IoU com o instante final fica entre 0,08 e 0,17
em Gliders e Spirals, e a média temporal fica vazia. Adotamos então a interpretação em que o frame
médio \emph{define o critério} e a máscara é gerada em cada frame:
\begin{equation}
  M(x,y,t) = \mathbb{1}\big[B(x,y,t) > \tau\big],\qquad \tau = 0{,}1507,
\end{equation}
em que $\tau$ é o limiar de Otsu \cite{otsu1979} do histograma de $B$ em todos os regimes do treino
($t>0$), que cai no vale entre o modo do fundo ($B\approx0$) e o do padrão ($B\approx0{,}25$ a
$0{,}45$). Um único $\tau$ para o fluido facilita a junção com a turma. Na validação, ele concorda
com $\tau$ por regime com IoU de 0,86 a 0,98, exceto em Gliders (0,74).

\paragraph{Vazamento de rótulo} Se $B$ fosse entrada, a rede só reaprenderia um limiar (IoU = 1).
Por isso \textbf{só $A$ entra na rede} e o rótulo vem de $B$: a rede precisa inferir a estrutura
de $B$ a partir de $A$. O problema é quase trivial em Bubbles e Spots, onde o melhor limiar único
em $A$ atinge IoU 0,97 a 0,98, e não trivial em Gliders e Spirals (0,81 a 0,85), sobretudo com um
limiar comum a todos os regimes.

\paragraph{Frames excluídos} $t=0$ (condições iniciais de Fourier com $A$ até 1,52 e $B$ negativo)
e os frames sem padrão ($\max B < 10^{-6}$), em que o padrão de Gliders e Spirals morreu e IoU e
Dice seriam indefinidos.

\subsection{Resolução}
\label{sec:resolucao}
A orientação foi aproximar a resolução de todos os fluidos da menor resolução usada no The Well,
mantendo a proporção e sem retirar as propriedades do fluido. O Gray-Scott já está em
$128\times128$, a menor resolução 2D do conjunto, e com proporção $1{:}1$. O fator de redução é,
portanto, 1, e \textbf{não há reamostragem}. Reduzir para $120\times120$ exigiria um fator não
inteiro ($0{,}9375$), com interpolação e filtragem \emph{antialiasing}, sobre estruturas de 4 a 8
pixels de largura. Pelo critério de Nyquist, isso borraria as interfaces e fundiria manchas
vizinhas, o que alteraria justamente a propriedade a segmentar. A nuvem de pontos tem $N=16\,384$
pontos por frame.

\subsection{Treinamento}
Cada amostra de treino é um bloco periódico $64\times64$ (4096 pontos) recortado em posição
aleatória do frame, como os blocos de cenas grandes no S3DIS \cite{armeni2016}. O aumento de dados
(rotações de $90^\circ$, reflexões, translações circulares e permutação aleatória dos pontos)
preserva a física, porque os campos são escalares e o domínio é periódico e isotrópico. Todas as
arquiteturas são treinadas por \textbf{20 épocas} (padronizado), com AdamW \cite{loshchilov2019}
(taxa máxima $10^{-3}$, \emph{weight decay} $10^{-4}$), política \emph{one-cycle} \cite{smith2019}
e perda $\mathcal{L} = \mathrm{BCE} + (1-\mathrm{Dice}_{\text{soft}})$ \cite{milletari2016}, que
compensa o desbalanceamento de Gliders (5\% de padrão). A semente é fixa (0), e o modelo avaliado é
o da época 20. Na inferência, o frame é coberto por 16 blocos $64\times64$ com passo 32 (cada ponto
é visto por 4 blocos), e as probabilidades são promediadas.

\subsection{Condições e seleção de parâmetros}
São 33 treinamentos (3 arquiteturas $\times$ 11 entradas: P0 e P1--P10). As probabilidades de
validação e teste são salvas, e os 11 pós-processamentos (Q0--Q10) são aplicados sem retreinar,
o que dá $3\times11\times11 = 363$ linhas de resultado: 3 de referência, 30 somente pré, 30 somente
pós e 300 pré e pós. Os parâmetros de cada Q são escolhidos na validação, e o teste é usado uma
única vez. Nas tabelas-resumo, a ``melhor'' configuração de cada condição também é escolhida pelo
Dice de validação.

\subsection{Formato padronizado das tabelas}
As tabelas seguem a ordem de colunas acordada com a turma: Fluido, Condição, Pré, Arquitetura,
Pós, IoU, Dice, Precision, Recall (média $\pm$ IC 95\%, 3 casas), Latência (ms) e Memória (MB).
As linhas seguem a ordem sequencial condição $\to$ arquitetura $\to$ pré $\to$ pós. O mesmo
esquema é salvo em \texttt{results/exp/metricas\_padrao.csv}, para a junção com os demais fluidos.

\section{Resultados}
\label{sec:resultados}

\ifthenelse{\equal{\resdir}{piloto}}{%
\noindent\fbox{\parbox{0.97\linewidth}{\small\textbf{Atenção:} esta versão mostra o \emph{piloto} reduzido
(CPU; 4 trajetórias de treino por regime, 4 épocas, pré P0, P1, P6 e P9, 3 trajetórias de teste por
regime e 4 blocos por frame). Ele valida o pipeline de ponta a ponta, mas \emph{não} é o experimento
do artigo. Os números finais vêm de \texttt{bash scripts/rodar\_tudo.sh} em GPU, com
\texttt{\textbackslash resdir} trocado para \texttt{exp}.}}\par\medskip}{}

\subsection{Tabelas visuais por instante}
As Figuras~\ref{fig:visual_gliders}--\ref{fig:visual_spots} mostram, para uma trajetória de teste de
cada regime, o instante inicial ($t=0$), dois instantes intermediários ($t=400$ s, na formação do
padrão, em que a variação entre frames é máxima, e $t=5000$ s, no meio da simulação) e o instante
final ($t=10\,000$ s). As linhas mostram o campo $A$ bruto, o campo após o pré-processamento mais
escolhido na condição pré e pós, o GT e, para cada arquitetura, o mapa de erros nas quatro condições
(VP em preto, FP em vermelho, FN em azul). A coluna $t=0$ é a condição inicial e fica fora da
avaliação quantitativa.

\newcommand{\visual}[2]{%
\begin{figure}[p]\centering
\includegraphics[height=0.92\textheight,width=\textwidth,keepaspectratio]{\figres{tabela_visual_#1}}
\caption{Tabela visual do regime #2: entrada, pré, GT e máscaras das três arquiteturas nas quatro
condições (pré/pós escolhidos na validação; legenda \emph{pré / pós}).}
\label{fig:visual_#1}
\end{figure}}
\visual{gliders}{Gliders}
\visual{bubbles}{Bubbles}
\visual{maze}{Maze}
\visual{worms}{Worms}
\visual{spirals}{Spirals}
\visual{spots}{Spots}

\subsection{Tabela de métricas padronizada}
A Tabela~\ref{tab:condicoes123} traz, em ordem sequencial, as condições sem pré e sem pós, somente
pré e somente pós. A condição pré e pós completa está no Apêndice (Tabela~\ref{tab:condicao4}), e a
Tabela~\ref{tab:resumo} resume, por arquitetura, a referência e a melhor configuração de cada
condição. A Figura~\ref{fig:heatmap} mostra as 363 combinações de uma vez: a célula
(nenhum, nenhum) é a referência, a linha ``nenhum'' é somente pós, a coluna ``nenhum'' é somente
pré e o restante é pré e pós.

\tabres{t_resumo}
\tabres{t_condicoes_123}

\begin{figure}[t]\centering
\includegraphics[width=\textwidth]{\figres{heatmap_dice}}
\caption{Dice no teste para cada par (pré, pós) e arquitetura. A cor é a diferença em relação à
referência (sem pré e sem pós, célula destacada).}
\label{fig:heatmap}
\end{figure}

\subsection{Resultados por regime, temporais e de custo}
\tabres{t_regimes}
\tabres{t_temporais}
\tabres{t_custo}

\begin{figure}[t]\centering
\includegraphics[width=\textwidth]{\figres{dice_por_regime}}
\caption{Dice no teste por regime (IC de 95\%) nas quatro condições, por arquitetura.}
\label{fig:regimes}
\end{figure}

\begin{figure}[t]\centering
\includegraphics[width=\textwidth]{\figres{curvas_treino}}
\caption{Perda de treino (contínua) e de validação (tracejada) por época, para cada pré-processamento.}
\label{fig:curvas}
\end{figure}

\section{Discussão: por que a segmentação melhora ou piora}
\label{sec:discussao}

A discussão parte de três fatos estruturais do problema, que explicam a maior parte dos efeitos:
\begin{enumerate}
  \item \textbf{Os dados não têm ruído.} A simulação é espectral e determinística
    (Seção~\ref{sec:grayscott}). Os filtros P6--P10 foram concebidos para remover ruído; aqui,
    qualquer alteração que produzem é alteração de \emph{sinal}.
  \item \textbf{O GT é um conjunto de nível de $B$ bruto, e $A$ é fortemente anticorrelacionado com
    $B$.} A fronteira correta fica próxima de um conjunto de nível de $A$, mas o nível ideal varia
    entre regimes ($A\approx0{,}45$ em Spirals a $0{,}72$ em Bubbles). Isso exige contexto, e não
    apenas um limiar.
  \item \textbf{As estruturas têm de 4 a 8 px de largura.} Operações com suporte de 3 px (filtros,
    elementos estruturantes) atuam na mesma escala do padrão e não são ``pequenas'' em relação a ele.
\end{enumerate}

\subsection{Arquiteturas}
O \textbf{PointNet} classifica cada ponto a partir de $(x_i,y_i,a_i)$ e do descritor global
$\mathbf{g}$. Na prática ele aprende um limiar sobre $a$, ajustado por $\mathbf{g}$. O descritor
global identifica implicitamente o regime (a distribuição de $a$ difere entre regimes), o que
permite deslocar o limiar por regime, algo que um limiar fixo não faz. A rede não tem acesso à
vizinhança, então não pode corrigir fronteiras deslocadas nem usar a forma das estruturas. Por isso
é a arquitetura mais sensível a pré-processamentos que deslocam conjuntos de nível (P6) e a que mais
pode ganhar com pós-processamentos que trazem contexto espacial (Q9, Q8).
O \textbf{PointNet++} agrega vizinhanças hierárquicas com coordenadas relativas e enxerga formas
de até dezenas de pixels nos níveis profundos. O \textbf{DGCNN} combina diferenças locais
($\mathbf{x}_j-\mathbf{x}_i$, análogas a um gradiente) com grafos dinâmicos que ligam pontos
semelhantes em todo o bloco. Espera-se que ambos superem o PointNet nas interfaces e nos regimes
não triviais (Gliders e Spirals). \revisar{Confirmar com a Tabela~\ref{tab:resumo} e a
Tabela~\ref{tab:regimes}: diferença de Dice na referência entre as redes, por regime.}

\subsection{Pré-processamentos numéricos (P1--P5)}
\paragraph{P1--P3 (afins)} Uma transformação $a' = \alpha a + \beta$ é absorvida pela primeira
camada linear: $w_a a' + b = (w_a\alpha)a + (w_a\beta + b)$. A normalização em lote
\cite{ioffe2015} logo em seguida remove o deslocamento e a escala por canal. Nas três redes, a
classe de funções representáveis é idêntica com e sem P1--P3. No DGCNN, o primeiro grafo usa só
$(x,y)$, e os grafos seguintes são calculados sobre características já normalizadas. O efeito
esperado é \emph{nulo dentro do IC}. Diferenças só podem vir da otimização: a escala relativa de
$a$ e das coordenadas muda a geometria da perda nos primeiros passos. \revisar{Verificar se
$|\Delta\mathrm{Dice}|$ de P1--P3 em relação a P0 fica dentro do IC de 95\%.} Esse é um resultado
negativo útil para o benchmark: para redes com BN, a normalização da entrada não é um fator
experimental relevante.

\paragraph{P4 (recorte)} Para $t>0$, $A\in[0{,}09;\,1{,}0006]$. O recorte em $Q_1$ e $Q_{99}$ só
achata o fundo dos vales (interior das manchas) e o platô, longe do nível em que fica a fronteira.
A fronteira é preservada, e o efeito esperado também é pequeno. Pode haver perda marginal onde a
profundidade do vale distingue estruturas recém-formadas, por exemplo na divisão de Spots.

\paragraph{P5 (log)} É monotônica e preserva os conjuntos de nível, então a mesma fronteira
continua representável. A transformação expande a faixa de $A$ baixo (interior) e comprime a faixa
perto de 1 (fundo). Isso aumenta a resolução numérica onde o limiar ideal é baixo (Spirals, Maze,
$A\approx0{,}45$ a $0{,}53$) e reduz a resolução onde ele é alto (Bubbles, $0{,}72$). O efeito
esperado é dependente do regime e de pequena magnitude. \revisar{Verificar na
Tabela~\ref{tab:regimes}.}

\subsection{Filtros espaciais (P6--P10)}
\paragraph{P6 (gaussiano)} A resposta em frequência é $\hat{G}(q) = e^{-\sigma^2 q^2/2}$. Para
estruturas com comprimento de onda $\lambda\approx8$ a $12$ px ($q = 2\pi/\lambda\approx0{,}5$ a
$0{,}8$ px$^{-1}$) e $\sigma = 1$ px, a amplitude cai para 73 a 88\%. O contraste entre vale e platô
diminui e os conjuntos de nível se deslocam: estruturas convexas e pequenas encolhem, e canais
estreitos se fecham. Como o GT é definido sobre $B$ \emph{sem} filtro, surge um viés sistemático
de fronteira que o PointNet não tem como corrigir. As redes com vizinhança podem compensar em parte,
porque sem ruído a deconvolução de uma atenuação moderada é bem posta. Espera-se que P6 piore,
mais no PointNet e nos regimes de estruturas pequenas (Gliders, Spots).

\paragraph{P7 (mediana $3\times3$)} Remove apenas estruturas com menos de $\sim$2 px e arredonda
cantos, preservando bordas retas. Para estruturas de 4 a 8 px, é quase a identidade, e o efeito
esperado é pequeno, negativo nas pontas de Worms e nos Gliders.

\paragraph{P8 (bilateral)} Com $\sigma_r = 0{,}05$, bem menor que o salto de $A$ na interface
($\sim$0,5), o peso de intensidade anula a mistura através da borda. A suavização ocorre só dentro
das regiões, que já são suaves. É praticamente a identidade, com efeito esperado nulo.

\paragraph{P9 (Perona-Malik)} O gradiente de $A$ nas interfaces ($\sim$0,1 por px) é maior que
$\kappa = 0{,}05$, de modo que $g(\|\nabla a\|)\approx e^{-4}$: a difusão para nas bordas e age
onde $\|\nabla a\|<\kappa$. O resultado é um campo levemente ``achatado'' por partes, com bordas
preservadas ou realçadas. Entre os filtros, é o único que poderia \emph{ajudar}, porque torna a
fronteira mais abrupta e reduz a ambiguidade do limiar. É também o mais caro (10 iterações).

\paragraph{P10 (wavelets)} O BayesShrink estima o ruído pela mediana dos coeficientes da subbanda
$HH_1$. Sem ruído, essa subbanda contém só detalhes finos do próprio sinal, concentrados nas
interfaces. O limiar resultante é pequeno, e o filtro é próximo da identidade, mas atenua
justamente os detalhes de borda que existem.

\paragraph{Síntese do pré} Os filtros espaciais trocam ruído (inexistente) por fidelidade de
fronteira. A previsão é que nenhum filtro melhore de forma significativa e que P6 seja o pior.
\revisar{Confirmar com a coluna ``nenhum'' da Figura~\ref{fig:heatmap} e com a seção ``Somente
pré'' da Tabela~\ref{tab:condicoes123}.} Em dados experimentais de reação-difusão (por exemplo,
\cite{lee1993}), com ruído de imagem real, a conclusão provavelmente se inverteria. Isso indica que
o efeito do pré-processamento depende da origem dos dados, e não só do método.

\subsection{Pós-processamentos (Q1--Q10)}
\paragraph{Q1 (limiar)} A perda BCE + Dice não produz probabilidades calibradas. O termo Dice
empurra as saídas para os extremos e desloca o ponto ótimo do limiar quando a classe positiva é
rara (Gliders, 5\%). Como $\tau^*$ é escolhido para maximizar o Dice de validação, Q1 tende a um
ganho pequeno e consistente, e funciona como \emph{piso} de comparação para os demais Q.

\paragraph{Q2--Q4 (morfologia)} A abertura remove FP isolados (salpicos), mas também Gliders e
pescoços de Spots em divisão, que têm a mesma escala do elemento estruturante. O fechamento une
estruturas vizinhas e fecha os corredores de fundo de Maze, que são estreitos (52\% do domínio é
padrão), o que gera FP. O efeito líquido depende do tipo de erro dominante da rede: abertura ajuda
se o erro dominante for FP disperso, fechamento ajuda se forem lacunas.
\revisar{Relacionar com Precision e Recall da Tabela~\ref{tab:condicoes123}.}

\paragraph{Q5--Q7 (componentes)} São seguros quando o parâmetro escolhido é pequeno ($a_{\min}$ de
2 a 4 px): eliminam FP espúrios sem tocar estruturas reais. Q6 tem um risco específico: em
\emph{Bubbles}, os ``buracos'' são as próprias bolhas de fundo, um objeto físico. Preenchê-los
apaga o padrão, e a validação tende a escolher $a_{\max}$ pequeno, o que limita o ganho. Q7
(probabilidade média do componente) funciona como uma histerese: componentes com pouca confiança
são removidos inteiros, em vez de pixel a pixel.

\paragraph{Q8 (watershed)} Desloca a fronteira do nível 0,5 para a crista de $\|\nabla P\|$.
Se a transição de $P$ for larga e simétrica, as duas coincidem e o efeito é nulo. Se for
assimétrica (típico com BN e Dice), a crista do gradiente aproxima melhor a interface.

\paragraph{Q9 (CRF)} É o único pós que reintroduz o \emph{dado} ($A$ bruto) no kernel de aparência.
Como a fronteira do GT acompanha um conjunto de nível de $A$, o termo bilateral alinha as bordas da
máscara às bordas de $A$ e corrige deslocamentos causados por um pré-processamento que borrou a
entrada (P6). Espera-se o maior ganho no PointNet (sem contexto local) e ganho menor no PointNet++ e
no DGCNN, que já usam vizinhança. \revisar{Verificar na Figura~\ref{fig:heatmap}, coluna Q9.}

\paragraph{Q10 (temporal)} Com frames a cada 200 s, a média móvel é inócua nos regimes estáticos
após o transiente (Maze, Spots) e reduz oscilações espúrias entre frames. Nos regimes móveis
(Gliders, Spirals), as estruturas se deslocam vários pixels entre frames, e a média borra a máscara
na direção do movimento. O efeito esperado é positivo nos estáticos e negativo nos móveis, e as
métricas temporais (Tabela~\ref{tab:temporais}) mostram esse contraste melhor que o Dice agregado.

\subsection{Interação pré $\times$ pós}
Os efeitos não são aditivos. Um pré que suaviza (P6) e um pós que suaviza (Q4, Q10) se somam na
direção de arredondar e unir estruturas, e o erro de fronteira se acumula. Já o CRF (Q9), por usar
$A$ \emph{bruto}, compensa parcialmente o borramento de P6. A Figura~\ref{fig:heatmap} permite
ler essas interações: uma combinação é sinérgica quando a célula supera a soma dos efeitos da sua
linha e da sua coluna.

\subsection{Evidências observadas}
\revisar{Os números abaixo são do \textbf{piloto} (4 épocas, 4 trajetórias de treino por regime,
pré P0, P1, P6 e P9, 3 trajetórias de teste por regime com frames a cada 1000~s). Substituir pelos
do experimento completo, mantendo a estrutura do argumento.}

\paragraph{Arquiteturas} Na referência, o DGCNN obteve Dice $0{,}948\pm0{,}011$, contra
$0{,}909\pm0{,}020$ do PointNet e $0{,}887\pm0{,}009$ do PointNet++. A vantagem do DGCNN aparece
sobretudo em Precision ($0{,}93$ contra $0{,}87$ e $0{,}83$): com diferenças locais
$\mathbf{x}_j-\mathbf{x}_i$, ele posiciona melhor as interfaces. O PointNet++ ficou abaixo do
PointNet no piloto, com a menor Precision. A interpolação 3-NN da decodificação suaviza as bordas
e produz FP na periferia das estruturas, e o modelo, mais profundo, é o mais prejudicado pelas
4 épocas. Gliders e Spirals foram os regimes mais difíceis para as três redes (Dice de 0,84 a 0,94),
como antecipava a análise de vazamento (Seção~\ref{sec:gt}).

\paragraph{Pré} O gaussiano (P6) foi o pior pré nas três redes ($\Delta$Dice de $-0{,}045$ no
PointNet, $-0{,}016$ no DGCNN e $-0{,}015$ no PointNet++). A perda foi maior no PointNet, que não
tem vizinhança para compensar o deslocamento das fronteiras, como previsto. Perona-Malik (P9) não
ajudou ($-0{,}015$, $-0{,}014$ e $-0{,}001$), o que confirma que, sem ruído, não há o que remover.
O Z-score (P1) foi neutro no PointNet ($0{,}909$ nos dois casos). Ele melhorou o PointNet++
($+0{,}011$) e o DGCNN ($+0{,}010$, no limite do IC). A invariância afim vale para a classe de
funções, mas não para a otimização em poucas épocas: com $a$ centrado e de variância unitária, as
primeiras atualizações ficam mais bem condicionadas. Resta verificar se o efeito persiste com
20 épocas.

\paragraph{Pós} Q1 deu ganho pequeno e consistente nas três redes ($+0{,}002$ a $+0{,}007$).
Q5--Q7 ficaram idênticos à referência: as redes não produzem componentes espúrios pequenos, então
não há o que remover. A abertura (Q2, Q4) piorou levemente o PointNet e o DGCNN, porque apagou
pontas e manchas pequenas. O CRF (Q9) foi o pós mais dependente da arquitetura: $+0{,}030$ no
PointNet++ (Precision de $0{,}835$ para $0{,}883$), por corrigir justamente as bordas suavizadas
pela interpolação, e leve piora no PointNet ($-0{,}004$) e no DGCNN ($-0{,}007$). A hipótese de que
o CRF ajudaria mais a rede sem vizinhança foi \emph{refutada}: o ganho depende do tipo de erro
(FP de borda), e não da ausência de contexto. A interação prevista entre P6 e Q9 apareceu no
PointNet: após o gaussiano, o CRF recupera parte da perda ($0{,}864\to0{,}878$), porque reintroduz
as bordas do $A$ bruto. O watershed (Q8) aumentou a Precision em todas as redes (por exemplo,
$0{,}868\to0{,}906$ no PointNet), com efeito pequeno no Dice.

\paragraph{Q10 e regimes móveis} A filtragem temporal foi o pós mais prejudicial: Dice de
$0{,}11$ a $0{,}13$ menor em todas as redes. A perda concentrou-se em Gliders (de $0{,}84$--$0{,}91$
para $\approx0{,}45$) e Spirals (de $0{,}86$--$0{,}94$ para $0{,}62$--$0{,}64$), enquanto em Maze e
Bubbles a variação ficou entre $0$ e $-0{,}015$. O erro de consistência temporal subiu de $0{,}008$ a $0{,}020$ para $0{,}10$ a
$0{,}12$. No piloto, o passo entre frames é de 1000 s, o que exagera o efeito. Mesmo assim, o
padrão confirma o mecanismo: a média temporal só é segura onde o padrão é estático.

\paragraph{Formação do padrão} Na tabela visual de Gliders (Figura~\ref{fig:visual_gliders}), em
$t=400$ s o GT é formado por \emph{anéis}: durante a formação, $B$ cai no centro das manchas.
As três redes preveem discos cheios, porque o vale de $A$ é cheio. Isso mostra o limite físico da
inferência de $B$ a partir de $A$ e indica por que o preenchimento de buracos (Q6) é perigoso
nesse instante.

\paragraph{Custo} Por frame, na CPU, a rede custa 322 ms (PointNet), 392 ms (PointNet++) e
1642 ms (DGCNN). Os pós morfológicos e de componentes custam de 1 a 2 ms, o watershed 23 ms e o
CRF 239 ms. O CRF quase dobra a latência do PointNet para um ganho que só aparece no PointNet++.

\subsection{Limitações}
\begin{itemize}
  \item O GT é um critério físico (limiar de $B$). A tarefa mede a capacidade de inferir a
    estrutura de $B$ a partir de $A$, e não de descobrir estruturas novas.
  \item Uma semente e um subconjunto de trajetórias, por custo. O IC por bootstrap captura a
    variação entre trajetórias, mas não a variação entre sementes.
  \item O treino em blocos $64\times64$ limita o contexto global a um quarto do domínio. Isso é
    suficiente para a escala das estruturas, mas não para a forma global de uma espiral.
  \item O CRF usa uma janela $7\times7$, e não conexões totais. O Q10 opera com passo de 200 s,
    grosso para regimes móveis.
\end{itemize}

\section{Conclusão}
\label{sec:conclusao}

Apresentamos a segmentação dos padrões de Turing do Gray-Scott segundo o protocolo comum da turma.
O GT é físico ($B>\tau$, congelado antes da comparação), e a entrada é restrita a $A$ para evitar
vazamento de rótulo. Usamos todas as 1200 trajetórias da divisão oficial, resolução nativa
$128\times128$ sem ampliação, 20 épocas com checkpoint da melhor época de validação, operações
periódicas e seleção de parâmetros só na validação. A análise mecanística prevê quatro resultados:
\begin{enumerate}
  \item normalizações afins (P1--P4) são neutras na U-Net, por causa da normalização em lote, mas
    podem afetar as redes pré-treinadas, que esperam entradas na escala do ImageNet;
  \item filtros espaciais (P6--P10) não ajudam em dados de simulação sem ruído, e o gaussiano piora
    a fronteira;
  \item a U-Net, que prediz na resolução completa, supera SegFormer e Mask2Former nas estruturas
    pequenas, cujos erros se concentram nas bordas reconstruídas a partir de $1/4$ da resolução;
  \item a filtragem temporal (Q5) só é segura nos regimes estáticos.
\end{enumerate}
No piloto, (2), (3) e (4) se confirmaram: nenhum filtro melhorou de forma significativa, a U-Net
superou os Transformers em 15 a 17 pontos de Dice, com erros deles concentrados em pontes de FP
entre estruturas vizinhas, e Q5 só foi inócuo em Maze e Bubbles. A hipótese (1) foi refinada: o
Z-score mudou a \emph{calibração}, e não a separabilidade. Com o limiar otimizado (Q1), ele produziu
a melhor configuração da U-Net (Dice 0,972). Por isso, no benchmark, o efeito de cada
pré-processamento deve ser lido junto com Q1.
\revisar{Atualizar com os números finais do experimento completo.} As tabelas estão no esquema
padronizado do protocolo, prontas para o artigo de benchmark entre sistemas físicos.


\section*{Disponibilidade de código e dados}
Código, configurações e resultados: repositório do projeto (\texttt{configs/experimento.yaml},
\texttt{scripts/04}--\texttt{08}, \texttt{PROTOCOLO.MD}). Os dados são lidos diretamente do repositório público do The Well
no Hugging Face (\texttt{polymathic-ai/gray\_scott\_reaction\_diffusion}).

\bibliographystyle{elsarticle-num}
\bibliography{referencias}

\appendix
\section{Resultados completos da condição pré e pós}
\tabres{t_condicao_4}

\end{document}
